% !TEX root = chapter-standalone.tex

\section{Sum of preorders}
\linkvideo{spring2021-tradeoffs:tradeoffs:orders:composing-posets:disj-poset} % Disjoint union of posets

Following the pattern, we can define the disjoint union of \SY{preorders} as a \SY{preorder} with the underlying set being the \SY{disjoint union} of the underlying sets.
\todotext{\bernina: after adding disjoint union of relations, then just use $\posAleq+\posBleq$}
\begin{ctdefinition}[Sum of preorders]
    \label{def:disjoint-union-of-posets}
    \SYNDEF{disjoint union of preorders}
    Given \SY{preorders}~$\posAdefinition$ and~$\posBdefinition$, their \emph{sum}~$\posA\Pplus\posB=\tupp{\posAset \setdisunion \posBset, {{\posleqof{\posA \Pplus \posB}}}}$, is the set $\posAset \setdisunion \posBset$ equipped with the relation~$\posleqof{\posA \Pplus \posB}$ defined by
    %
    \begin{equation} \label{eq:poset-disunion}
        \begin{aligned}
            \tup{i, \posela} \posleqof{\posA\Pplus\posB}  \tup{ j, \poselb}
             & \Leftrightarrow &
            \begin{cases}
                i = j, \text{ and}                            \\
                \posela \posAleq \poselb \text{ if } i = 1, \\
                \posela \posBleq \poselb \text{ if } i = 2.
            \end{cases}
        \end{aligned}
    \end{equation}
    %        \begin{equation} \label{eq:poset-disunion}
    %        \begin{aligned}
    %            \posleqof{\posA\Pplus\posB}\colon
    %            (\posAset\setdisunion\posBset)\cartprod (\posAset
    %            \setdisunion\posBset)                           & \sto \boolset, \\
    %            \tupp{\disunionA{\posela},\disunionA{\poselb} } &
    %            \mapsto (\posela\posAleq \poselb), \\
    %            \tupp{\disunionB{\cdot},\disunionA{\cdot} }     &
    %            \mapsto \false, \\
    %            \tupp{\disunionA{\cdot},\disunionB{\cdot}}      &
    %            \mapsto \false, \\
    %            \tupp{\disunionB{\posela},\disunionB{\poselb}}  &
    %            \mapsto (\posela\posBleq \poselb).
    %        \end{aligned}
    %    \end{equation}
\end{ctdefinition}

\begin{lemma}
    \label{lem:sum-preorder}
    Let~\posA and~\posB be \SY{preorders}.
    Then the relation~$\posleqof{\posA\Pplus\posB}$ is reflexive and transitive, so~$\posA\Pplus\posB$ is a \SY{preorder}.
    If, moreover,~\posA and~\posB are \SY{posets}, then~$\posA\Pplus\posB$ is a \SY{poset}.
\end{lemma}
\begin{proof}
    In each case, related elements lie in the same summand, so each property reduces to the corresponding property of~$\posAleq$ or of~$\posBleq$.
    \begin{enumerate}
        \item Reflexivity: let~$\tup{i, \posela}\setin \posAset \setdisunion \posBset$.
        Since~$i = i$, and~$\posela \posAleq \posela$ (if~$i=1$) or~$\posela \posBleq \posela$ (if~$i=2$) by reflexivity, we have~$\tup{i, \posela} \posleqof{\posA\Pplus\posB} \tup{i, \posela}$.
        \item Transitivity: suppose~$\tup{i, \posela} \posleqof{\posA\Pplus\posB} \tup{j, \poselb}$ and~$\tup{j, \poselb} \posleqof{\posA\Pplus\posB} \tup{k, \poselc}$.
        By definition,~$i = j$ and~$j = k$, so all three elements lie in the same summand.
        If~$i = 1$, then~$\posela \posAleq \poselb$ and~$\poselb \posAleq \poselc$, so~$\posela \posAleq \poselc$ by transitivity of~$\posAleq$; if~$i = 2$, the same argument works with~$\posBleq$.
        In either case,~$\tup{i, \posela} \posleqof{\posA\Pplus\posB} \tup{k, \poselc}$.
        \item Antisymmetry (when~\posA and~\posB are \SY{posets}): suppose~$\tup{i, \posela} \posleqof{\posA\Pplus\posB} \tup{j, \poselb}$ and~$\tup{j, \poselb} \posleqof{\posA\Pplus\posB} \tup{i, \posela}$.
        Then~$i = j$, and~$\posela$ and~$\poselb$ are related in both directions in the same summand, so~$\posela = \poselb$ by antisymmetry of~$\posAleq$ (if~$i=1$) or of~$\posBleq$ (if~$i=2$).
        Hence~$\tup{i, \posela} = \tup{j, \poselb}$.
    \end{enumerate}
\end{proof}

The expression \cref{eq:poset-disunion} can be intimidating at first, but all it is saying is that the relation of the sum is obtained by stitching together the two relations~$\posAleq$ and~$\posBleq$.
No element of $\posAset$ is related to an element of $\posBset$, and vice versa.

\begin{example}
    Consider the \SY{posets}~$\posA,\posB$, over the sets~$\posAset=\makeset{\sbretzel, \sfondue}$ with~$\sbretzel \posAleq \sfondue$, and~$\posBset=\makeset{\stea,\smilk}$, with~$\smilk \posBleq \stea$.
    Their \SY{disjoint union} can be represented as in \cref{fig:poset-coproduct}.
    %
    \begin{figure}[h!]
        \centering
        \includesag{40_disjoint_union}
        \caption{Disjoint union of \SY{posets}.}
        \label{fig:poset-coproduct}
    \end{figure}
\end{example}
\vfill
\begin{gradedexercise}[\exname{MeasurePosetSum}]
    \label{ex:MeasurePosetSum}
    Let~\posA and~\posB be \SY{posets}.
    Prove the following properties:
    \begin{enumerate}
        \item The width of the sum is the sum of the widths:
              \begin{equation}\label{eq:MeasurePosetSum-1}
                  \posetwidth(\posA \Pplus \posB) = \posetwidth(\posA) + \posetwidth(\posB).
              \end{equation}
        \item The height of the sum is the maximum of the heights:
              \begin{equation}\label{eq:MeasurePosetSum-2}
                  \posetheight(\posA \Pplus \posB) = \max(\posetheight(\posA), \posetheight(\posB)).
              \end{equation}
    \end{enumerate}
\end{gradedexercise}
\solutionof{MeasurePosetSum}
